% ============================================================================
\section{Introduction}
\label{sec:intro}

This section covers the quantitative tools from the introduction lecture: how to
measure time and throughput, the CPU equation, instruction types and CPI, and
Amdahl's Law.

% ----------------------------------------------------------------------------
\subsection{What to know}

\subsubsection{Time and throughput}
\begin{definition}[Latency and throughput]
\textbf{Latency} (execution time, response time) is the time to finish
\emph{one} task. \textbf{Throughput} (bandwidth) is the number of tasks
finished per unit time.
\end{definition}
\begin{itemize}[leftmargin=*]
  \item If a machine runs one task at a time, throughput $=1/\text{latency}$.
  With overlap (several tasks in flight: pipelining, multiple cores, more
  passengers per plane) throughput can rise while latency stays the same or
  even gets worse.
  \item \textbf{Elapsed (wall-clock) time} includes I/O and other programs.
  \textbf{CPU time} $=$ user CPU time $+$ system CPU time; it excludes waiting.
  \item Unix \texttt{time}: \texttt{90.7u 12.9s 2:39 65\%} means
  $(90.7+12.9)/159=65\%$ of the elapsed time was CPU time. With several
  threads on several cores the percentage can exceed 100\%.
\end{itemize}

\subsubsection{Performance and speedup}
\[
  \text{Performance}=\frac{1}{\text{ExTime}}\qquad
  \text{Speedup of A over B}=\frac{\text{Perf}_A}{\text{Perf}_B}
  =\frac{\text{ExTime}_B}{\text{ExTime}_A}=\frac{\text{old time}}{\text{new time}}
\]
``X is $n\%$ faster than Y'' means
$\dfrac{\text{ExTime}_Y}{\text{ExTime}_X}=1+\dfrac{n}{100}$, so
$n=100\,\dfrac{\text{ExTime}_Y-\text{ExTime}_X}{\text{ExTime}_X}$ (divide by
the time of the \emph{faster} machine).

\subsubsection{The CPU equation}
\begin{theorem}[CPU time]
\[
  T=\frac{\text{Instructions}}{\text{Program}}\times
    \frac{\text{Cycles}}{\text{Instruction}}\times
    \frac{\text{Seconds}}{\text{Cycle}}
   =I\times\CPI\times C=\frac{I\times\CPI}{f}.
\]
\end{theorem}
\begin{center}
\begin{tabular}{@{}lccc@{}}
\toprule
Affects $\dots$ & $I$ & $\CPI$ & $C$ \\
\midrule
Program, compiler, ISA & X & X & \\
Organization (CPU design) & & X & X \\
Technology (VLSI) & & & X \\
\bottomrule
\end{tabular}
\end{center}

\subsubsection{Instruction types and CPI}
\begin{prop}[CPI of an instruction mix]
With $I_i$ instructions of class $i$, each taking $\CPI_i$ cycles, and
frequency $F_i=I_i/I$:
\[
  \text{Cycles}=\sum_i \CPI_i\,I_i,\qquad
  \CPI=\frac{\text{Cycles}}{I}=\sum_i \CPI_i\,F_i,
\]
\[
  \text{fraction of time in class } i=\frac{\CPI_i\,F_i}{\CPI}.
\]
\end{prop}
\begin{remark}
CPI is an \emph{average}, so it depends on the mix. A change that removes cheap
instructions can \emph{raise} the CPI while lowering the time: always compare
total cycles or time, never CPI or instruction count on their own.
\end{remark}

\subsubsection{Amdahl's Law}
\begin{theorem}[Amdahl's Law]
If an enhancement speeds up a fraction $F$ of the \emph{original execution
time} by a factor $S$, and the rest is unchanged,
\[
  \text{ExTime}_{new}=\text{ExTime}_{old}\left[(1-F)+\frac FS\right],\qquad
  \text{Speedup}=\frac{1}{(1-F)+\dfrac FS}\;\le\;\frac{1}{1-F}.
\]
\end{theorem}
\begin{itemize}[leftmargin=*]
  \item \textbf{Several disjoint enhancements} (every $F_i$ a fraction of the
  original time):
  $\text{Speedup}=\dfrac{1}{(1-\sum F_i)+\sum F_i/S_i}$.
  \item \textbf{Parallel processors}: with parallel fraction $p$ on $N$
  processors, $\text{Speedup}=\dfrac{1}{(1-p)+p/N}\le\dfrac1{1-p}$.
  \item \textbf{Target speedup $T$}: $S=\dfrac{F}{1/T-(1-F)}$; impossible when
  $T\ge1/(1-F)$. Required fraction: $F=\dfrac{1-1/T}{1-1/S}$.
  \item \textbf{$F$ from an instruction mix}: $F=\CPI_i F_i/\CPI$, the share of
  \emph{time}, not of instructions.
  \item \textbf{$f'$ measured after the enhancement} (share of the \emph{new}
  time spent in enhanced mode): $\text{Speedup}=(1-f')+f'S$, and the original
  fraction is $F=f'S/\big((1-f')+f'S\big)$.
  \item \textbf{Enhancement that slows the clock} by a factor $k$:
  $\text{Speedup}=\dfrac{1}{\big[(1-F)+F/S\big]\,k}$.
\end{itemize}

% ----------------------------------------------------------------------------
\subsection{Common traps}
\begin{enumerate}[leftmargin=*]
  \item Speedup is old time over new time, never the reverse; a speedup below 1
  is a slowdown.
  \item ``$n\%$ faster'' divides by the faster time: 10~s vs.\ 8~s is 25\%
  faster, while 8~s vs.\ 10~s is 20\% \emph{less time}.
  \item Converting clock rate: 2.5~GHz is a 0.4~ns cycle. Keep $f$ in Hz and
  $T$ in seconds until the last step.
  \item A higher clock rate or a lower CPI on its own does not mean faster; check
  all three factors of $I\times\CPI\times C$.
  \item In Amdahl's Law, $F$ is a fraction of \emph{time}. Given a mix, convert
  with $\CPI_iF_i/\CPI$ first.
  \item A fraction measured on the \emph{enhanced} machine is not the $F$ of the
  formula; use $(1-f')+f'S$.
  \item Cross-check every Amdahl answer with the CPU equation when you can:
  same $I$ and clock $\Rightarrow$ speedup $=\CPI_{old}/\CPI_{new}$.
\end{enumerate}

% ----------------------------------------------------------------------------
\subsection{Lecturer's review questions}
From the lecturer's review slides (\texttt{Review\_final\_2026.pdf}), with worked answers.

\begin{exercise}[Lecturer 1]\label{l:perf1}
In a program's execution time, multiplication instructions take 60\% and division
instructions 30\%. Multiplication can be made up to $6\times$ faster and division up
to $12\times$ faster. To make the computer $2\times$ faster by improving only
\emph{one} component, which component is selected?
(a) multiplication; (b) division; (c) both a and b are correct; (d) both a and b are
wrong.
\end{exercise}
\begin{solution}
\textbf{(a)}. Multiplication: $\dfrac{1}{0.4+0.6/6}=\dfrac{1}{0.5}=2.0$, which
reaches the target. Division: $\dfrac{1}{0.7+0.3/12}=\dfrac{1}{0.725}=1.38$; even
$S\to\infty$ gives only $1/0.7=1.43<2$.
\end{solution}

\begin{exercise}[Lecturer 2]
If both multiplication and division are improved, the new computer is faster by:
(1) 4.000; (2) 4.175; (3) 4.444; (4) none of the above.
\end{exercise}
\begin{solution}
\textbf{(3)}: $\dfrac{1}{(1-0.6-0.3)+0.6/6+0.3/12}=\dfrac{1}{0.1+0.1+0.025}
=\dfrac{1}{0.225}=4.444$.
\end{solution}

\begin{exercise}[Lecturer 3]
A program of 4 million instructions runs on a 2.5~GHz CPU in 8~ms. The average CPI
is: (a) 5; (b) 10; (c) 20; (d) none of the above.
\end{exercise}
\begin{solution}
\textbf{(a)}: $\CPI=\dfrac{T\times f}{I}=\dfrac{8\times10^{-3}\times2.5\times10^9}
{4\times10^6}=\dfrac{2\times10^7}{4\times10^6}=5$.
\end{solution}

\begin{exercise}[Lecturer 4]
A new compiler produces 5 million instructions with CPI 2; the clock is unchanged.
The speedup is: (a) 1.5; (b) 4; (c) 2; (d) none of the above.
\end{exercise}
\begin{solution}
\textbf{(c)}: $T'=\dfrac{5\times10^6\times2}{2.5\times10^9}=4$~ms, speedup
$=8/4=2$. More instructions, but far fewer cycles ($10^7$ vs.\ $2\times10^7$).
\end{solution}

% ----------------------------------------------------------------------------
\subsection{Exam-style questions}
Solutions are in Section~\ref{sec:intro-sol}.

\begin{exercise}[Latency vs.\ throughput]\label{q:intro-lat}
Web server A handles one request at a time and takes 2~ms per request. Web
server B handles 4 requests at the same time, each taking 5~ms.
\begin{enumerate}[label=(\alph*),leftmargin=*]
  \item Which server has the lower latency, and by what factor?
  \item Find the throughput of each server in requests per second. Which is
  higher, and by what percentage?
  \item 3000 requests arrive together. How long does each server take to finish
  all of them?
  \item Which server would you pick for one interactive user, and which for a
  busy data centre? Why?
  \item A program run on B reports \texttt{12.0u 2.0s 0:07 200\%}. Explain how
  the CPU percentage can exceed 100\%.
\end{enumerate}
\end{exercise}

\begin{exercise}[Three design options]\label{q:intro-design}
A program executes $4\times10^9$ instructions with $\CPI=1.6$ on a 2~GHz
processor. The design team has two options:
\begin{itemize}[leftmargin=*]
  \item \textbf{Option A}: a deeper pipeline raises the clock to 2.5~GHz, but
  the CPI rises to 1.9.
  \item \textbf{Option B}: new instructions cut the instruction count by 15\%;
  the clock stays at 2~GHz but the CPI rises to 1.75.
\end{itemize}
\begin{enumerate}[label=(\alph*),leftmargin=*]
  \item Find the CPU time of the original design and of each option.
  \item Which option is better? Give its speedup over the original and say how
  many percent faster it is.
  \item What CPI would option A need to match option B?
  \item Which factors of the CPU equation does each option change, and which
  part of the design (compiler, ISA, organization, technology) is responsible?
\end{enumerate}
\end{exercise}

\begin{exercise}[Instruction classes]\label{q:intro-mix}
A 2~GHz processor runs a program of $5\times10^9$ instructions with this mix:
\begin{center}
\begin{tabular}{@{}lcc@{}}
\toprule
Class & Frequency & CPI \\
\midrule
ALU & 40\% & 1 \\
Load/Store & 30\% & 3 \\
Branch & 20\% & 2 \\
Floating point & 10\% & 6 \\
\bottomrule
\end{tabular}
\end{center}
\begin{enumerate}[label=(\alph*),leftmargin=*]
  \item Find the average CPI.
  \item Find the fraction of execution time spent in each class.
  \item Find the CPU time and the MIPS rating.
  \item A compiler optimization removes half of the branches and changes nothing
  else. Find the new instruction count, CPI and CPU time. Explain why the CPI
  goes \emph{up} although the program gets faster.
\end{enumerate}
\end{exercise}

\begin{exercise}[Choosing an enhancement]\label{q:intro-choose}
Use the original mix of Exercise~\ref{q:intro-mix}. Two improvements are
proposed:
\begin{itemize}[leftmargin=*]
  \item \textbf{E1}: a faster FP unit lowers the FP CPI from 6 to 2.
  \item \textbf{E2}: a better cache lowers the load/store CPI from 3 to 2.
\end{itemize}
\begin{enumerate}[label=(\alph*),leftmargin=*]
  \item For each, find $F$ and $S$ and the overall speedup with Amdahl's Law.
  \item Check both answers with the CPU equation.
  \item Only 10\% of the instructions are FP but 30\% are loads/stores. Explain
  why E1 still wins.
  \item Find the speedup with both enhancements, with Amdahl's Law and with the
  CPU equation.
\end{enumerate}
\end{exercise}

\begin{exercise}[Limits and targets]\label{q:intro-limit}
A program spends 60\% of its execution time in matrix multiplication.
\begin{enumerate}[label=(\alph*),leftmargin=*]
  \item An accelerator runs matrix multiplication $20\times$ faster. Find the
  overall speedup.
  \item What is the largest overall speedup any accelerator could give?
  \item How fast must the accelerator be for an overall speedup of 2.4?
  \item Can an overall speedup of 2.5 be reached? Explain.
  \item What fraction of the time would have to be accelerated by $20\times$ to
  reach an overall speedup of 4?
\end{enumerate}
\end{exercise}

\begin{exercise}[Fraction measured after the enhancement]\label{q:intro-after}
After a SIMD unit is added (it runs vectorizable code $10\times$ faster), a
profiler on the \emph{new} machine shows that 40\% of the execution time is
spent in SIMD code.
\begin{enumerate}[label=(\alph*),leftmargin=*]
  \item Find the overall speedup due to the SIMD unit.
  \item What fraction of the \emph{original} execution time was vectorizable?
  \item A student plugs $F=0.4$, $S=10$ into Amdahl's Law. What does the
  student get, and why is it wrong?
  \item The SIMD unit is improved to $20\times$. Find the speedup over the
  original machine (without SIMD).
\end{enumerate}
\end{exercise}

\begin{exercise}[Enhancement with a clock penalty]\label{q:intro-clock}
An enhancement speeds up 30\% of the execution time by $4\times$, but its extra
logic makes the clock cycle 5\% longer.
\begin{enumerate}[label=(\alph*),leftmargin=*]
  \item Find the overall speedup with and without the clock penalty.
  \item How much longer could the clock cycle become before the enhancement
  stops paying off?
\end{enumerate}
\end{exercise}

\begin{exercise}[Amdahl's Law for multiprocessors]\label{q:intro-par}
90\% of a program can run in parallel.
\begin{enumerate}[label=(\alph*),leftmargin=*]
  \item Find the speedup on 8 and on 64 processors.
  \item What is the speedup limit?
  \item How many processors are needed for a speedup of 8?
  \item Comment on the efficiency (speedup per processor) in (a).
\end{enumerate}
\end{exercise}

\begin{exercise}[True or false]\label{q:intro-tf}
Answer and give a one-line reason.
\begin{enumerate}[label=(\alph*),leftmargin=*]
  \item A processor with a higher clock rate is always faster.
  \item Of two code sequences for the same task, the one with the lower CPI is
  always faster.
  \item On a machine that runs one task at a time, halving the latency doubles
  the throughput.
  \item In Amdahl's Law, $F$ can be taken as the fraction of instructions that
  benefit.
  \item With enough money, any overall speedup can be reached by speeding up
  only part of a program.
  \item If X is 25\% faster than Y, then Y takes 25\% more time than X.
\end{enumerate}
\end{exercise}

% ----------------------------------------------------------------------------
\subsection{Solutions}
\label{sec:intro-sol}

\begin{noot}{Exercise~\ref{q:intro-lat}}
\begin{enumerate}[label=(\alph*),leftmargin=*]
  \item A: 2~ms, B: 5~ms. A has $5/2=2.5\times$ lower latency.
  \item A: $1/0.002=500$ requests/s. B: $4/0.005=800$ requests/s. B has
  $800/500=1.6\times$, i.e.\ 60\% higher throughput.
  \item A: $3000\times2\text{ ms}=6$~s. B: $3000/4=750$ rounds of 5~ms
  $=3.75$~s.
  \item One interactive user sees only latency, so A. A data centre with many
  users cares about requests served per second, so B. Latency and throughput can
  rank machines in opposite orders.
  \item CPU time $=12.0+2.0=14$~s over an elapsed time of 7~s gives
  $14/7=200\%$. CPU time is summed over all cores, so two threads running on two
  cores at the same time give 2~s of CPU time per second of wall-clock time.
\end{enumerate}
\end{noot}

\begin{noot}{Exercise~\ref{q:intro-design}}
\begin{enumerate}[label=(\alph*),leftmargin=*]
  \item Original: $T=\dfrac{4\times10^9\times1.6}{2\times10^9}=3.2$~s.\\
  A: $T_A=\dfrac{4\times10^9\times1.9}{2.5\times10^9}=3.04$~s.\\
  B: $I_B=0.85\times4\times10^9=3.4\times10^9$,
  $T_B=\dfrac{3.4\times10^9\times1.75}{2\times10^9}=2.975$~s.
  \item B is better. Speedup $=3.2/2.975=1.076$, so B is 7.6\% faster than the
  original (A gives $3.2/3.04=1.053$). B is also
  $100\,(3.04-2.975)/2.975=2.2\%$ faster than A.
  \item $\CPI_A=\dfrac{T_B\times f_A}{I}=\dfrac{2.975\times2.5\times10^9}{4\times10^9}=1.859$.
  \item A changes $C$ (shorter cycle) and CPI: organization. B changes $I$
  and CPI: ISA (and the compiler that uses the new instructions).
\end{enumerate}
\end{noot}

\begin{noot}{Exercise~\ref{q:intro-mix}}
\begin{enumerate}[label=(\alph*),leftmargin=*]
  \item $\CPI=0.4(1)+0.3(3)+0.2(2)+0.1(6)=0.4+0.9+0.4+0.6=2.3$.
  \item ALU $0.4/2.3=17.4\%$, load/store $0.9/2.3=39.1\%$, branch
  $0.4/2.3=17.4\%$, FP $0.6/2.3=26.1\%$.
  \item $T=\dfrac{5\times10^9\times2.3}{2\times10^9}=5.75$~s.
  $\text{MIPS}=\dfrac{2\times10^9}{2.3\times10^6}=869.6$.
  \item Branches drop from $1\times10^9$ to $0.5\times10^9$, so
  $I'=4.5\times10^9$. Cycles $=11.5\times10^9-0.5\times10^9\times2=10.5\times10^9$.
  $\CPI'=10.5/4.5=2.33$, $T'=10.5\times10^9/(2\times10^9)=5.25$~s, a speedup of
  $5.75/5.25=1.095$. The removed branches (CPI 2) were cheaper than the average
  (2.3), so the remaining mix is more expensive per instruction, but there are
  fewer cycles in total. (MIPS also falls, to 857, although the program is
  faster.)
\end{enumerate}
\end{noot}

\begin{noot}{Exercise~\ref{q:intro-choose}}
\begin{enumerate}[label=(\alph*),leftmargin=*]
  \item E1: $F=0.6/2.3=0.261$, $S=6/2=3$:
  $\text{Speedup}=\dfrac{1}{0.739+0.261/3}=\dfrac{1}{0.739+0.087}=1.21$.\\
  E2: $F=0.9/2.3=0.391$, $S=3/2=1.5$:
  $\text{Speedup}=\dfrac{1}{0.609+0.391/1.5}=\dfrac{1}{0.609+0.261}=1.15$.
  \item E1: $\CPI'=2.3-0.1(6-2)=1.9$, $2.3/1.9=1.21$.
  E2: $\CPI'=2.3-0.3(3-2)=2.0$, $2.3/2.0=1.15$. Both match.
  \item What matters is the time saved per instruction,
  $F_i\,(\CPI_{old}-\CPI_{new})$: E1 saves $0.1\times4=0.4$ cycles, E2 saves
  $0.3\times1=0.3$ cycles. FP instructions are rare but expensive, so they take
  26\% of the time.
  \item Amdahl: $\dfrac{1}{(1-0.261-0.391)+0.261/3+0.391/1.5}
  =\dfrac{1}{0.348+0.087+0.261}=\dfrac{1}{0.696}=1.44$.
  CPU equation: $\CPI'=2.3-0.4-0.3=1.6$, $2.3/1.6=1.44$.
\end{enumerate}
\end{noot}

\begin{noot}{Exercise~\ref{q:intro-limit}}
\begin{enumerate}[label=(\alph*),leftmargin=*]
  \item $\dfrac{1}{0.4+0.6/20}=\dfrac{1}{0.43}=2.33$.
  \item $S\to\infty$: $1/(1-0.6)=2.5$.
  \item $S=\dfrac{0.6}{1/2.4-0.4}=\dfrac{0.6}{0.01667}=36$.
  \item No: 2.5 is the limit. It would need $1/2.5=0.4=(1-F)+F/S$, i.e.\
  $F/S=0$, an infinitely fast accelerator.
  \item $F=\dfrac{1-1/4}{1-1/20}=\dfrac{0.75}{0.95}=0.789$: about 79\% of the time.
\end{enumerate}
\end{noot}

\begin{noot}{Exercise~\ref{q:intro-after}}
\begin{enumerate}[label=(\alph*),leftmargin=*]
  \item Take the new time as 1: 0.6 is non-SIMD and 0.4 is SIMD. Without the
  unit the SIMD part would take $0.4\times10=4$, so the old time is
  $0.6+4=4.6$. Speedup $=4.6$.
  \item $F=4/4.6=0.870$: 87\% of the original time.
  \item $1/(0.6+0.04)=1.56$. The 40\% is a share of the \emph{new}, shorter
  time; Amdahl's $F$ must be a share of the \emph{old} time. Check:
  $1/(0.130+0.870/10)=1/0.217=4.6$.
  \item $F=20/23$: $\dfrac{1}{3/23+(20/23)/20}=\dfrac{1}{4/23}=5.75$.
\end{enumerate}
\end{noot}

\begin{noot}{Exercise~\ref{q:intro-clock}}
\begin{enumerate}[label=(\alph*),leftmargin=*]
  \item Without the penalty: new time $=0.7+0.3/4=0.775$, speedup
  $1/0.775=1.29$. With it: $0.775\times1.05=0.814$, speedup $1.23$.
  \item It pays off while $0.775\,k<1$, i.e.\ $k<1/0.775=1.29$: the cycle can
  grow by up to about 29\%.
\end{enumerate}
\end{noot}

\begin{noot}{Exercise~\ref{q:intro-par}}
\begin{enumerate}[label=(\alph*),leftmargin=*]
  \item $N=8$: $\dfrac{1}{0.1+0.9/8}=\dfrac{1}{0.2125}=4.71$.
  $N=64$: $\dfrac{1}{0.1+0.9/64}=\dfrac{1}{0.1141}=8.77$.
  \item $1/0.1=10$, however many processors.
  \item $1/8=0.125=0.1+0.9/N\Rightarrow N=0.9/0.025=36$.
  \item Efficiency: $4.71/8=59\%$ on 8 processors but only $8.77/64=14\%$ on 64.
  The serial 10\% dominates as $N$ grows.
\end{enumerate}
\end{noot}

\begin{noot}{Exercise~\ref{q:intro-tf}}
\begin{enumerate}[label=(\alph*),leftmargin=*]
  \item False: time is $I\times\CPI/f$; a higher CPI (Exercise~\ref{q:intro-design},
  option A) or a larger $I$ can cancel the faster clock.
  \item False: the lower-CPI sequence may execute more instructions; compare
  total cycles.
  \item True: throughput $=1/\text{latency}$ when tasks do not overlap.
  \item False: $F$ is a fraction of execution time; convert with
  $\CPI_iF_i/\CPI$.
  \item False: the speedup is bounded by $1/(1-F)$.
  \item True: $\text{ExTime}_Y/\text{ExTime}_X=1.25$. (But X takes 20\% less
  time than Y.)
\end{enumerate}
\end{noot}

